% Options for packages loaded elsewhere
% Options for packages loaded elsewhere
\PassOptionsToPackage{unicode}{hyperref}
\PassOptionsToPackage{hyphens}{url}
\PassOptionsToPackage{dvipsnames,svgnames,x11names}{xcolor}
%
\documentclass[
  letterpaper,
  DIV=11,
  numbers=noendperiod]{scrartcl}
\usepackage{xcolor}
\usepackage[top=2.0cm,bottom=2.0cm,left=2.5cm,right=2.5cm,footskip=0.7cm]{geometry}
\usepackage{amsmath,amssymb}
\setcounter{secnumdepth}{-\maxdimen} % remove section numbering
\usepackage{iftex}
\ifPDFTeX
  \usepackage[T1]{fontenc}
  \usepackage[utf8]{inputenc}
  \usepackage{textcomp} % provide euro and other symbols
\else % if luatex or xetex
  \usepackage{unicode-math} % this also loads fontspec
  \defaultfontfeatures{Scale=MatchLowercase}
  \defaultfontfeatures[\rmfamily]{Ligatures=TeX,Scale=1}
\fi
\usepackage{lmodern}
\ifPDFTeX\else
  % xetex/luatex font selection
\fi
% Use upquote if available, for straight quotes in verbatim environments
\IfFileExists{upquote.sty}{\usepackage{upquote}}{}
\IfFileExists{microtype.sty}{% use microtype if available
  \usepackage[]{microtype}
  \UseMicrotypeSet[protrusion]{basicmath} % disable protrusion for tt fonts
}{}
\makeatletter
\@ifundefined{KOMAClassName}{% if non-KOMA class
  \IfFileExists{parskip.sty}{%
    \usepackage{parskip}
  }{% else
    \setlength{\parindent}{0pt}
    \setlength{\parskip}{6pt plus 2pt minus 1pt}}
}{% if KOMA class
  \KOMAoptions{parskip=half}}
\makeatother
% Make \paragraph and \subparagraph free-standing
\makeatletter
\ifx\paragraph\undefined\else
  \let\oldparagraph\paragraph
  \renewcommand{\paragraph}{
    \@ifstar
      \xxxParagraphStar
      \xxxParagraphNoStar
  }
  \newcommand{\xxxParagraphStar}[1]{\oldparagraph*{#1}\mbox{}}
  \newcommand{\xxxParagraphNoStar}[1]{\oldparagraph{#1}\mbox{}}
\fi
\ifx\subparagraph\undefined\else
  \let\oldsubparagraph\subparagraph
  \renewcommand{\subparagraph}{
    \@ifstar
      \xxxSubParagraphStar
      \xxxSubParagraphNoStar
  }
  \newcommand{\xxxSubParagraphStar}[1]{\oldsubparagraph*{#1}\mbox{}}
  \newcommand{\xxxSubParagraphNoStar}[1]{\oldsubparagraph{#1}\mbox{}}
\fi
\makeatother


\usepackage{longtable,booktabs,array}
\newcounter{none} % for unnumbered tables
\usepackage{calc} % for calculating minipage widths
% Correct order of tables after \paragraph or \subparagraph
\usepackage{etoolbox}
\makeatletter
\patchcmd\longtable{\par}{\if@noskipsec\mbox{}\fi\par}{}{}
\makeatother
% Allow footnotes in longtable head/foot
\IfFileExists{footnotehyper.sty}{\usepackage{footnotehyper}}{\usepackage{footnote}}
\makesavenoteenv{longtable}
\usepackage{graphicx}
\makeatletter
\newsavebox\pandoc@box
\newcommand*\pandocbounded[1]{% scales image to fit in text height/width
  \sbox\pandoc@box{#1}%
  \Gscale@div\@tempa{\textheight}{\dimexpr\ht\pandoc@box+\dp\pandoc@box\relax}%
  \Gscale@div\@tempb{\linewidth}{\wd\pandoc@box}%
  \ifdim\@tempb\p@<\@tempa\p@\let\@tempa\@tempb\fi% select the smaller of both
  \ifdim\@tempa\p@<\p@\scalebox{\@tempa}{\usebox\pandoc@box}%
  \else\usebox{\pandoc@box}%
  \fi%
}
% Set default figure placement to htbp
\def\fps@figure{htbp}
\makeatother





\setlength{\emergencystretch}{3em} % prevent overfull lines

\providecommand{\tightlist}{%
  \setlength{\itemsep}{0pt}\setlength{\parskip}{0pt}}



 


\KOMAoption{captions}{tableheading}
\makeatletter
\@ifpackageloaded{caption}{}{\usepackage{caption}}
\AtBeginDocument{%
\ifdefined\contentsname
  \renewcommand*\contentsname{Table of contents}
\else
  \newcommand\contentsname{Table of contents}
\fi
\ifdefined\listfigurename
  \renewcommand*\listfigurename{List of Figures}
\else
  \newcommand\listfigurename{List of Figures}
\fi
\ifdefined\listtablename
  \renewcommand*\listtablename{List of Tables}
\else
  \newcommand\listtablename{List of Tables}
\fi
\ifdefined\figurename
  \renewcommand*\figurename{Figure}
\else
  \newcommand\figurename{Figure}
\fi
\ifdefined\tablename
  \renewcommand*\tablename{Table}
\else
  \newcommand\tablename{Table}
\fi
}
\@ifpackageloaded{float}{}{\usepackage{float}}
\floatstyle{ruled}
\@ifundefined{c@chapter}{\newfloat{codelisting}{h}{lop}}{\newfloat{codelisting}{h}{lop}[chapter]}
\floatname{codelisting}{Listing}
\newcommand*\listoflistings{\listof{codelisting}{List of Listings}}
\makeatother
\makeatletter
\makeatother
\makeatletter
\@ifpackageloaded{caption}{}{\usepackage{caption}}
\@ifpackageloaded{subcaption}{}{\usepackage{subcaption}}
\makeatother
\usepackage{bookmark}
\IfFileExists{xurl.sty}{\usepackage{xurl}}{} % add URL line breaks if available
\urlstyle{same}
% fallback for those not using the hyperref driver hyperxmp:
\makeatletter
\@ifundefined{xmpquote}{\newcommand{\xmpquote}[1]{#1}}{}
\makeatother
\hypersetup{
  pdftitle={TTex Employee Fraud Analysis},
  colorlinks=true,
  linkcolor={blue},
  filecolor={Maroon},
  citecolor={Blue},
  urlcolor={Blue},
  pdfcreator={LaTeX via pandoc}}


\title{TTex Employee Fraud Analysis}
\author{}
\date{}
\begin{document}
\maketitle

\renewcommand*\contentsname{Table of Contents}
{
\setcounter{tocdepth}{3}
\tableofcontents
}

\subsection{1. Executive Summary}\label{executive-summary}

This analysis focuses on analyzing company vehicle usage, trip
distances, and fuel consumption to identify potential irregularities and
assess data quality. Data from 64 files was cleaned, validated, and
merged into a single dataset. The analysis included verifying employee
and vehicle activity periods, examining trip distances, and evaluating
fuel consumption patterns.

The analysis identified 30 trips with potentially unusual distances and
significant fuel consumption discrepancies, particularly involving two
employees from the Transport Department using a Fiat Ducato and a Ford
Transit.

The findings highlight potential data quality issues and possible
overreporting of fuel consumption that require further investigation.

\subsection{2. Problem And Objectives}\label{problem-and-objectives}

The company needs to ensure the accuracy of its vehicle usage records
and identify potential irregularities in reported trip distances and
fuel consumption.

The main objectives of the analysis were to: - Trip distance analysis:
Identify unusual distances between the same locations and detect
potentially inaccurate trip records. - Fuel consumption analysis:
Evaluate fuel efficiency across vehicles and identify unusually high
reported fuel consumption. - Employee activity verification: Ensure that
trips were made by active employees using active company vehicles. -
Risk identification: Highlight potentially suspicious trips, vehicles,
and employees requiring further investigation.

The ultimate goal is to identify potential operational inefficiencies,
and support better control of company vehicle usage and fuel expenses.

\subsection{3. Data Validatiion And
Peparation}\label{data-validatiion-and-peparation}

Data from all weeks (64 files in total) was combined into five separate
DataFrames: df\_department, df\_position, df\_employee, df\_car and
df\_trip. Duplicates resulting from the concatenation of the files were
removed. The correctness of the data types was verified. After data
cleaning, the suffix ``\_clean'' was added to the names of all
DataFrames.

\begin{enumerate}
\def\labelenumi{\arabic{enumi}.}
\tightlist
\item
  The \textbf{df\_department} DataFrame contains 5 unique records, each
  representing an organizational department and consisting of three
  fields: id, name, and description.
\item
  The \textbf{df\_position} DataFrame contains 36 unique records, the
  IDs and names of job positions within the company.
\item
  The \textbf{df\_employee} DataFrame contains 36 unique records with
  employee information, including first names, last names, gender,
  department IDs, position IDs, dates of birth, hire dates, and
  employment termination dates.
\item
  The \textbf{df\_car} DataFrame contains 20 unique records related to
  vehicles used within the company. It includes the following columns:
  car ID, brand, model, year of manufacture, fuel type, transport usage,
  activity status, activity start date, and activity end date. Missing
  values were identified in the activity status column and filled with
  appropriate values based on the activity dates.
\item
  The \textbf{df\_trip} DataFrame contains trip IDs, start and end
  dates, car IDs, employee IDs, distance traveled, fuel consumption, and
  start and end locations. The DataFrame contained no duplicate records.
  Incorrect values in the from and to columns were corrected. After data
  cleaning, 13 unique locations remained in these columns, including
  cities and other places, such as offices and shops.
\end{enumerate}

Finally, all cleaned DataFrames were merged into a single DataFrame
named df\_total, as shown below.

\begin{longtable}[]{@{}ll@{}}
\caption{Total DataFrame Structure}\label{T_56666}\tabularnewline
\toprule\noalign{}
column & dtype \\
\midrule\noalign{}
\endfirsthead
\toprule\noalign{}
column & dtype \\
\midrule\noalign{}
\endhead
\bottomrule\noalign{}
\endlastfoot
trip\_id & int64 \\
trip\_start & datetime64{[}us{]} \\
trip\_end & datetime64{[}us{]} \\
trip\_distance & int64 \\
trip\_fuel & float64 \\
trip\_from & str \\
trip\_to & str \\
car\_brand & str \\
car\_model & str \\
car\_year & int64 \\
car\_fuel\_type & str \\
car\_transport & bool \\
car\_active & str \\
car\_active\_from & datetime64{[}us{]} \\
car\_active\_to & datetime64{[}us{]} \\
employee\_name & str \\
employee\_surname & str \\
employee\_gender & str \\
employee\_year\_of\_birth & int64 \\
employee\_date\_of\_employment & datetime64{[}us{]} \\
employee\_employment\_end\_date & datetime64{[}us{]} \\
employee\_position\_name & str \\
employee\_department\_name & str \\
employee\_department\_description & str \\
\end{longtable}

\begin{verbatim}
Shape: (4432, 24)
\end{verbatim}

\subsection{4. Exploratory Data
Analysis}\label{exploratory-data-analysis}

A total of 443 trips were completed. Trip \textbf{distances} ranged from
3 km (between Olsztyn and the shop or office) to 576 km (between
Szczecin and Olsztyn), with an average distance of 150.34 km.

\begin{figure}[H]

{\centering \pandocbounded{\includegraphics[keepaspectratio]{fraud_analysis_files/figure-pdf/cell-10-output-1.pdf}}

}

\caption{Distribution of Trip Distances}

\end{figure}%

The \textbf{trip dates were validated and found to be consistent}. In
all cases, the trip start date was earlier than or equal to the end
date.

No \textbf{trips involving inactive vehicles were detected}, i.e., trips
that started after the vehicle's activity end date or started before
that date but ended after it.

No trips made by employees before their employment start date or after
their employment end date were detected.

Most intercity routes show a very similar relative range of distances,
approximately 7--8\%. \vspace{1cm}

\begin{longtable}[]{@{}lllllllll@{}}
\caption{Descriptive Statistics of Trip Distances by
Route}\label{T_25378}\tabularnewline
\toprule\noalign{}
route & count & mean & std & min & 25\% & 50\% & 75\% & max \\
\midrule\noalign{}
\endfirsthead
\toprule\noalign{}
route & count & mean & std & min & 25\% & 50\% & 75\% & max \\
\midrule\noalign{}
\endhead
\bottomrule\noalign{}
\endlastfoot
Bartoszyce ↔ Olsztyn & 536.00 & 73.53 & 1.61 & 71.00 & 72.00 & 73.50 &
75.00 & 76.00 \\
Bydgoszcz ↔ Olsztyn & 544.00 & 219.54 & 4.60 & 212.00 & 215.00 & 220.00
& 223.00 & 228.00 \\
Elbląg ↔ Olsztyn & 378.00 & 113.92 & 2.44 & 110.00 & 112.00 & 114.00 &
116.00 & 118.00 \\
Gdańsk ↔ Olsztyn & 486.00 & 197.39 & 4.02 & 190.00 & 194.00 & 197.00 &
201.00 & 204.00 \\
Malbork ↔ Olsztyn & 494.00 & 134.87 & 2.77 & 130.00 & 133.00 & 135.00 &
137.00 & 140.00 \\
Mikołajki ↔ Olsztyn & 570.00 & 87.09 & 1.82 & 84.00 & 85.00 & 87.00 &
89.00 & 90.00 \\
Olsztyn ↔ Ostróda & 320.00 & 41.52 & 0.90 & 40.00 & 41.00 & 42.00 &
42.00 & 43.00 \\
Olsztyn ↔ Sklep & 38.00 & 4.58 & 1.62 & 3.00 & 3.00 & 4.00 & 5.00 &
8.00 \\
Olsztyn ↔ Szczecin & 90.00 & 557.58 & 11.86 & 538.00 & 547.00 & 560.00 &
566.00 & 576.00 \\
Olsztyn ↔ Toruń & 510.00 & 186.36 & 3.78 & 180.00 & 183.00 & 186.00 &
189.00 & 193.00 \\
Olsztyn ↔ Urząd & 58.00 & 6.41 & 2.36 & 4.00 & 4.00 & 7.00 & 7.00 &
11.00 \\
Olsztyn ↔ Warszawa & 408.00 & 228.27 & 4.70 & 220.00 & 224.00 & 228.50 &
232.00 & 236.00 \\
\end{longtable}

The median distance was calculated for each route, and trips exceeding
the median distance by more than 10\% were identified. A total of 30
potentially suspicious trips were identified.

\begin{longtable}[]{@{}lllllll@{}}
\caption{Suspicious Trips}\label{T_40626}\tabularnewline
\toprule\noalign{}
route & distance & Mdn & Mdn\_dev & E\_name & E\_surname & count \\
\midrule\noalign{}
\endfirsthead
\toprule\noalign{}
route & distance & Mdn & Mdn\_dev & E\_name & E\_surname & count \\
\midrule\noalign{}
\endhead
\bottomrule\noalign{}
\endlastfoot
Olsztyn ↔ Sklep & 8 & 4.00 & 100.00 & Mariola & Suchań & 2 \\
Olsztyn ↔ Sklep & 7 & 4.00 & 75.00 & Fryderyk & Przybylski & 2 \\
Olsztyn ↔ Sklep & 7 & 4.00 & 75.00 & Jadwiga & Gall & 4 \\
Olsztyn ↔ Urząd & 11 & 7.00 & 57.14 & Ernest & Rząca & 2 \\
Olsztyn ↔ Urząd & 11 & 7.00 & 57.14 & Roksana & Czachor & 2 \\
Olsztyn ↔ Urząd & 10 & 7.00 & 42.86 & Roksana & Czachor & 2 \\
Olsztyn ↔ Urząd & 10 & 7.00 & 42.86 & Fryderyk & Przybylski & 2 \\
Olsztyn ↔ Urząd & 10 & 7.00 & 42.86 & Juliusz & Ikonowicz & 4 \\
Olsztyn ↔ Sklep & 5 & 4.00 & 25.00 & Fryderyk & Przybylski & 2 \\
Olsztyn ↔ Sklep & 5 & 4.00 & 25.00 & Emil & Sekula & 6 \\
Olsztyn ↔ Sklep & 5 & 4.00 & 25.00 & Ernest & Rząca & 2 \\
\end{longtable}

\scriptsize Legend: Mdn - median, Mdn\_dev - distance deviation pct, E -
employee.\\

\vspace{0.5cm}
\normalsize

The \textbf{average fuel consumption per 100 km} was calculated for each
vehicle used. The minimum and maximum values differ significantly from
the mean and median, which may indicate substantial inaccuracies in the
reported fuel consumption data. \textbf{A threshold of 50\% above the
median fuel consumption for each vehicle} was adopted as an indicator of
potentially incorrect fuel consumption values for individual trips. As a
result, \textbf{670 trips with potential irregularities were
identified}.

\begin{longtable}[]{@{}llllll@{}}
\caption{Fuel Consumption per 100 km per
car}\label{T_d6f3b}\tabularnewline
\toprule\noalign{}
count & mean & median & std & min & max \\
\midrule\noalign{}
\endfirsthead
\toprule\noalign{}
count & mean & median & std & min & max \\
\midrule\noalign{}
\endhead
\bottomrule\noalign{}
\endlastfoot
1210 & 11.05 & 9.53 & 4.55 & 7.87 & 27.44 \\
614 & 11.02 & 9.14 & 4.52 & 7.87 & 27.44 \\
1126 & 12.57 & 10.22 & 5.68 & 8.80 & 30.71 \\
1196 & 12.49 & 10.65 & 5.34 & 8.67 & 30.71 \\
286 & 7.57 & 7.65 & 0.39 & 5.42 & 7.75 \\
\end{longtable}

Two employees were identified as significantly overreporting fuel
consumption. These employees traveled in a Fiat Ducato and a Ford
Transit on nearly all routes. In the table, Avg\_dev\_\% represents the
average percentage deviation in fuel consumption from the median,
calculated across all trips made by each employee, including those
without detected irregularities.

\begin{longtable}[]{@{}lllll@{}}
\caption{Suspicious Employees}\label{T_0fae9}\tabularnewline
\toprule\noalign{}
E\_name & E\_surname & department & suspicious\_trips & Avg\_dev\_\% \\
\midrule\noalign{}
\endfirsthead
\toprule\noalign{}
E\_name & E\_surname & department & suspicious\_trips & Avg\_dev\_\% \\
\midrule\noalign{}
\endhead
\bottomrule\noalign{}
\endlastfoot
Magdalena & Dobrzyńska & Transport & 544 & 102.77 \\
Emil & Sekula & Transport & 126 & 30.78 \\
\end{longtable}

\subsection{5. Recomendation}\label{recomendation}

First, discrepancies in trip distances that deviated significantly from
the median should be investigated. Most of these involved trips to
offices and shops, where the destinations were not precisely defined.
Therefore, these differences do not necessarily indicate irregularities.

The significantly higher fuel consumption reported by employees in the
Transport Department also requires further investigation. Fuel
consumption levels nearly twice the expected values strongly suggest
that the reported amount of fuel used may have been substantially
overestimated.




\end{document}
